% ID: SEL-D0011-BC-E05
% ESTADO: cotejada. Segunda pasada visual completa y borrador compilado revisados.
% Fuente: Selectividad/matematicas_aplicadas/EXTRA_SUPLENTE_Modelo-1-Enunciados (Mat-Ap-2021).pdf
% SHA-256: 4a6c7d34ff35c3e7678c52970b4d4558cee2c35224b46c66035cb0bc8d5a7ec8
% Página: 2; bloque C; ejercicio original 5.

Un equipo andaluz de baloncesto jugó en una temporada un 40\% de los partidos
en casa y el resto fuera. De los partidos que jugó en casa, obtuvo un 60\% de
victorias y el resto fueron derrotas, mientras que de los que jugó fuera,
obtuvo un 30\% de victorias y el resto derrotas. Se elige un partido de este
equipo al azar.
\begin{enumerate}[label=\alph*)]
  \item \textbf{(1.25 puntos)} Calcule la probabilidad de que el partido
  acabase en victoria.
  \item \textbf{(1 punto)} Calcule la probabilidad de que el partido haya
  sido jugado en casa, sabiendo que el resultado final fue una derrota.
  \item \textbf{(0.25 puntos)} Si además se sabe que el 10\% de las victorias
  obtenidas en casa y el 20\% de las obtenidas fuera se produjeron tras una
  prórroga, calcule la probabilidad de que el partido acabase en victoria
  y que además esa victoria haya sido tras una prórroga.
\end{enumerate}
